% ---------------------------------------------------------------
%  Muc II.5 -- Kha nang ap dung va hieu qua
% ---------------------------------------------------------------
\subsection{Khả năng áp dụng và hiệu quả của sáng kiến}

\subsubsection{Khả năng áp dụng}

\textbf{Tại đơn vị.} Sáng kiến đã được áp dụng tại \DONVI\ cho cả giáo viên tổ
Toán và học sinh các lớp \dotfill\ trong năm học \NAMHOC.

\textbf{Khả năng nhân rộng.} Do phần khung của hệ thống tách rời hoàn toàn khỏi
nội dung câu hỏi, việc nhân rộng chỉ là bổ sung dữ liệu, không phải viết lại
phần mềm:

\begin{itemize}[leftmargin=1.2cm, itemsep=3pt]
  \item \textbf{Sang khối lớp khác:} bổ sung phân phối chương trình, bảng yêu
        cầu cần đạt và các hàm sinh câu hỏi của lớp 11, lớp 12. Toàn bộ thuật
        toán dựng ma trận, chọn câu, sinh đề, chấm bài dùng lại nguyên vẹn.
  \item \textbf{Sang trường khác:} mỗi trường chỉ cần một tài khoản giáo viên
        kết nối Google Classroom của mình; hệ thống chạy trên web nên không
        phải cài đặt gì.
  \item \textbf{Sang môn học khác:} các môn có câu hỏi dạng tính toán (Vật lí,
        Hoá học, Sinh học) dùng lại được toàn bộ phần mềm, chỉ cần viết hàm
        sinh câu hỏi của môn mình theo đúng chuẩn mã định danh ở Giải pháp 1.
\end{itemize}

\textbf{Điều kiện áp dụng.} Máy tính hoặc điện thoại có kết nối Internet; giáo
viên biết dùng máy tính ở mức cơ bản; lớp học Google Classroom (không bắt
buộc). Riêng việc \textit{bổ sung câu hỏi mới} vào ngân hàng đòi hỏi giáo viên
biết viết hàm Python cơ bản --- đây là ngưỡng duy nhất cần vượt qua, và theo
kinh nghiệm của tôi thì một giáo viên Toán làm quen được sau vài buổi, vì bản
chất chỉ là viết lại công thức toán mà mình vốn đã thuộc.

\subsubsection{Hiệu quả đối với giáo viên}

\begin{center}
\begin{tabular}{|p{6.4cm}|c|c|c|}
\hline
\textbf{Nội dung} & \textbf{Trước khi áp dụng} & \textbf{Sau khi áp dụng} &
\textbf{Chênh lệch} \\ \hline
Thời gian soạn một đề kiểm tra định kì & & & \\ \hline
Thời gian soạn thêm một mã đề & & & \\ \hline
Thời gian làm đáp án, biểu điểm & & & \\ \hline
Thời gian chấm một lớp (phần trắc nghiệm) & & & \\ \hline
Số mã đề thực sự khác nhau dùng được & & & \\ \hline
\end{tabular}
\end{center}
\textit{Bảng 7. Hiệu quả về thời gian làm việc của giáo viên}

\ghichu{Cách lấy số cho bảng này: nhờ hai đến ba đồng nghiệp bấm giờ một lần
soạn đề theo cách cũ, rồi bấm giờ một lần dùng hệ thống. Chỉ cần số phút, không
cần khảo sát lớn. Riêng dòng ``thời gian soạn thêm một mã đề'' là dòng thuyết
phục nhất --- theo cách cũ gần như phải soạn lại cả đề, còn ở đây là bấm thêm
một lần.}

\subsubsection{Hiệu quả đối với học sinh}

\begin{center}
\begin{tabular}{|c|p{7.6cm}|c|c|}
\hline
\textbf{TT} & \textbf{Nội dung khảo sát sau khi áp dụng} & \textbf{Số HS} &
\textbf{Tỉ lệ} \\ \hline
1 & Thấy hứng thú hơn khi được tự tạo đề luyện tập & & \\ \hline
2 & Cho rằng biết điểm ngay giúp mình nhớ bài hơn & & \\ \hline
3 & Có luyện lại bằng nút ``Làm đề khác cùng cấu trúc'' & & \\ \hline
4 & Cho rằng kiểm tra công bằng hơn vì mỗi bạn một đề & & \\ \hline
5 & Tự đánh giá là tiến bộ hơn so với đầu năm & & \\ \hline
\end{tabular}
\end{center}
\textit{Bảng 8. Kết quả khảo sát học sinh sau khi áp dụng (tổng số: \dots\ học sinh)}

\begin{center}
\begin{tabular}{|p{3.4cm}|c|c|c|c|c|}
\hline
\multirow{2}{*}{\textbf{Lớp}} & \multirow{2}{*}{\textbf{Sĩ số}} &
\multicolumn{4}{c|}{\textbf{Xếp loại điểm bài kiểm tra}} \\ \cline{3-6}
 & & \textbf{Giỏi} & \textbf{Khá} & \textbf{Trung bình} & \textbf{Yếu} \\ \hline
Lớp thực nghiệm \dotfill & & & & & \\ \hline
Lớp đối chứng \dotfill & & & & & \\ \hline
\end{tabular}
\end{center}
\textit{Bảng 9. So sánh kết quả bài kiểm tra giữa lớp thực nghiệm và lớp đối chứng}

\ghichu{Bảng 9 là bảng hội đồng chấm đọc kĩ nhất. Lưu ý ba điều: (1) hai lớp
phải tương đương nhau về đầu vào --- nếu không thì ghi rõ điểm trung bình môn
Toán học kì trước của hai lớp để chứng minh; (2) dùng cùng một đề hoặc hai đề
cùng ma trận; (3) ghi thêm điểm trung bình của mỗi lớp dưới bảng. Nếu chị chưa
có lớp đối chứng, có thể thay bằng so sánh cùng một lớp trước và sau khi áp
dụng --- vẫn hợp lệ, chỉ cần nói rõ.}

\subsubsection{Hiệu quả đo được từ chính hệ thống}

Ngoài số liệu khảo sát, bản thân hệ thống ghi lại nhật kí hoạt động nên có thể
lấy ra những con số khách quan sau:

\begin{center}
\begin{tabular}{|p{8.4cm}|c|}
\hline
\textbf{Chỉ số} & \textbf{Số liệu} \\ \hline
Số hàm sinh câu hỏi đã viết trong ngân hàng & \\ \hline
Số mã định danh câu hỏi đã có hàm sinh & \\ \hline
Số đề đã được sinh ra trong năm học & \\ \hline
Số lượt học sinh làm bài trực tuyến & \\ \hline
Số lớp đã kết nối Google Classroom & \\ \hline
\end{tabular}
\end{center}
\textit{Bảng 10. Số liệu vận hành thực tế của hệ thống}

\ghichu{Hai dòng đầu lấy bằng cách đếm trong thư mục ngân hàng Python (tại thời
điểm viết mục này là 46 hàm sinh ứng với 32 mã định danh, đều thuộc chương I
lớp 10). Ba dòng sau lấy bằng cách đếm số dòng trong các bảng
\texttt{de\_da\_sinh}, \texttt{exam\_history} và \texttt{classroom\_coursework}
của cơ sở dữ liệu --- nhớ lấy số vào sát ngày nộp, và ghi rõ mốc thời gian
thống kê ngay dưới bảng.}

\subsubsection{Hiệu quả về kinh tế và xã hội}

\textbf{Về kinh tế.} Hệ thống tự xây dựng nên không phát sinh chi phí mua phần
mềm tạo đề thương mại hằng năm. Chi phí duy trì chỉ gồm tiền thuê máy chủ và
tên miền, ở mức một cá nhân chi trả được. Việc chấm phần trắc nghiệm không gọi
trí tuệ nhân tạo nên không phát sinh chi phí theo số lượt làm bài, dù có bao
nhiêu học sinh cùng sử dụng. Ngoài ra, học sinh không phải mua thêm sách bài
tập để luyện, vì hệ thống sinh ra bài luyện không giới hạn.

\textbf{Về xã hội.} Thời gian giáo viên tiết kiệm được ở khâu soạn đề và chấm
bài được chuyển sang khâu có giá trị hơn: phân tích chỗ hổng của học sinh và
dạy lại. Học sinh ở mọi điều kiện gia đình đều tiếp cận được nguồn đề đúng
chương trình và có lời giải đúng, không phụ thuộc vào việc có đi học thêm hay
có mua được sách tham khảo hay không --- đây là ý nghĩa mà tôi tâm đắc nhất khi
làm sáng kiến này.

\clearpage
